\scanpage{39}
% archive foliation: 38
Passage à la limite

\uline{\unsure{Thm}} Soient $I$ un ordonné \deleted{filtrant}, $X,Y$ dans $(\mathrm{Cat})$, $(X_i)_{i\in I}$ et $Y=(Y_i)_{i\in I}$ deux filtrations par des sous-cat. fermées (resp. ouvertes), \deleted{\unsure{$f$}} $X=\bigcup X_i$, $Y=\bigcup Y_i$.
\[
f:X\longrightarrow Y
\]
tel que $f(X_i)\subset Y_i$, soit $f_i:X_i\unsure{\mathbin{-}|\longrightarrow}Y_i$.
Supposons que les $f_i$ soient des équiv. faibles, alors $f$ est une équiv. faible.

\margin{\unsure{On suppose que les inf existent et que les filtrations} \unsure{commutent}, \unsure{dans le cas où} $i=\operatorname{Inf}j_\alpha$, \unsure{on ait} $X_i=\bigcap_\alpha X_{j_\alpha}$.}

On est ramené par dualité au cas des ss-catégories ouvertes. On va \deleted{\unsure{supposer}} interpréter la filtration de $X$ comme la donnée de
\diagram{\begin{tikzcd}X\arrow[r,"\varphi"]&I\end{tikzcd}}
défini par
\[
\varphi(x)=\operatorname{Inf}_{\substack{i\text{ dans }I(x)=\\\{i\in I\mid x\in X_i\}}}i
\]
--- on suppose donc que dans $I$ les Inf \inserted{(des parties non vides)} existent. On \deleted{\unsure{pourra}} \inserted{aura}, pour
\[
\alpha:x\longrightarrow y
\]
$\varphi(x)\leq\varphi(y)$ (car $I(y)\subset I(x)$, donc $\operatorname{Inf}I(y)\geq\operatorname{Inf}I(x)$) d'où la valeur de $\varphi(\alpha)$. On \deleted{\unsure{aura supposé, où}} \inserted{devra supposer}
\[
X_i=\varphi^{-1}(I_i)\qquad\text{où}\qquad I_i=\{j\in I\mid j\leq i\}
\]
\deleted{\unsure{$\varphi$}} i.e. $\underbrace{x\in X_i}_{\text{i.e. }i\in I(x)}\Longleftrightarrow\varphi(x)\leq i$

\margin{\uline{NB} $\varphi(x)$ est ainsi le plus petit élément de $I(x)$ et $I(x)=\{j\in I\mid j\geq\varphi(x)\}$.}
